\documentclass[11pt]{article}
\usepackage{mathblatt}
\usepackage{adjustbox,varwidth}
% gesetzt von werkzeuge/setzer.py; Lösungen nach bau/bauregeln.md „Lösungen“.
\newcommand{\kennung}{MZE}
\begin{document}
\pfheftstil
\fancyfoot[C]{{\footnotesize\color{mbgrau}\kennung\hfill\thepage}}
\pfheftkopf{Termwerte berechnen}{Klasse 7 \textperiodcentered{} Lösungen}

\begin{pfloesung}{}
\lz{1b)}{$6;\ 11;\ 46$}{a) $5 \cdot 2 - 4 = 10 - 4 = 6$; b) $5 \cdot 3 - 4 = 15 - 4 = 11$; c) $5 \cdot 10 - 4 = 50 - 4 = 46$}{}
\end{pfloesung}
\begin{pfloesung}{}
\lz{2.}{$13;\ 20;\ 11;\ 8$}{a) $7 + 2 \cdot 3 = 7 + 6 = 13$; b) $4 \cdot (3 + 2) = 4 \cdot 5 = 20$; c) $20 - 3 \cdot 3 = 20 - 9 = 11$; d) $2 \cdot 3 + 3 - 1 = 6 + 3 - 1 = 8$}{}
\end{pfloesung}
\begin{pfloesung}{}
\lz{3.}{$1;\ 4;\ 7;\ 10;\ 13$}{$3 \cdot 1 - 2 = 1$; $3 \cdot 2 - 2 = 4$; $3 \cdot 3 - 2 = 7$; $3 \cdot 4 - 2 = 10$; $3 \cdot 5 - 2 = 13$\par\smallskip \adjustbox{max width=\linewidth,max height=4cm}{\begin{varwidth}{50cm}\wertetabelle[1,4,7,10,13]{x}{3x - 2}{1,2,3,4,5}\end{varwidth}}}{}
\end{pfloesung}
\begin{pfloesung}{}
\lz{4.}{nein, $22$\,€; es fehlen $2$\,€}{$4 + 3 \cdot 6 = 4 + 18 = 22$\,€. Das Geld reicht nicht, es fehlen $22 - 20 = 2$\,€.}{}
\end{pfloesung}
\begin{pfloesung}{}
\lz{5b)}{$4;\ -5;\ 7$}{a) $3 \cdot (-2) + 10 = -6 + 10 = 4$; b) $3 \cdot (-5) + 10 = -15 + 10 = -5$; c) $3 \cdot (-1) + 10 = -3 + 10 = 7$}{}
\end{pfloesung}
\begin{pfloesung}{}
\lz{6.}{$11;\ 6;\ 17;\ -9$}{a) $8 - (-3) = 8 + 3 = 11$; b) $-2 \cdot (-3) = 6$; c) $5 - 4 \cdot (-3) = 5 - (-12) = 5 + 12 = 17$; d) $-3 - 6 = -9$}{}
\end{pfloesung}
\begin{pfloesung}{}
\lz{7.}{$2;\ 18$}{a) $6 \cdot (-2 - 1) + 20 = 6 \cdot (-3) + 20 = -18 + 20 = 2$; b) $10 - 2 \cdot (-7 + 3) = 10 - 2 \cdot (-4) = 10 - (-8) = 18$}{}
\end{pfloesung}
\begin{pfloesung}{}
\lz{8.}{$3 - 4x$ (Wert $11$), um $3$ größer}{$3 - 4 \cdot (-2) = 3 + 8 = 11$; $2 \cdot (-2 + 6) = 2 \cdot 4 = 8$. $3 - 4x$ hat den größeren Wert, um $11 - 8 = 3$.}{}
\end{pfloesung}
\begin{pfloesung}{}
\lz{9b)}{$25;\ 25;\ 1;\ 100$}{a) $5^2 = 25$; b) $(-5)^2 = (-5) \cdot (-5) = 25$; c) $(-1)^2 = 1$; d) $(-10)^2 = 100$}{}
\end{pfloesung}
\begin{pfloesung}{}
\lz{10.}{$16;\ 8;\ 6$}{a) $4 \cdot (-2)^2 = 4 \cdot 4 = 16$; b) $(-2)^2 + 4 = 4 + 4 = 8$; c) $(-2)^2 - (-2) = 4 + 2 = 6$}{}
\end{pfloesung}
\begin{pfloesung}{}
\lz{11.}{$1;\ -3$}{a) $10 - (-3)^2 = 10 - 9 = 1$; b) $2 \cdot (-1)^2 + 5 \cdot (-1) = 2 \cdot 1 + (-5) = 2 - 5 = -3$}{}
\end{pfloesung}
\begin{pfloesung}{}
\lz{12.}{$8;\ 3;\ 0;\ -1;\ 0;\ 3$; negativ nur für $x = 1$}{$(-2)^2 - 2 \cdot (-2) = 4 + 4 = 8$; $(-1)^2 - 2 \cdot (-1) = 1 + 2 = 3$; $0^2 - 2 \cdot 0 = 0$; $1^2 - 2 \cdot 1 = 1 - 2 = -1$; $2^2 - 2 \cdot 2 = 0$; $3^2 - 2 \cdot 3 = 9 - 6 = 3$. Negativ nur für $x = 1$.\par\smallskip \adjustbox{max width=\linewidth,max height=4cm}{\begin{varwidth}{50cm}\wertetabelle[8,3,0,-1,0,3]{x}{x^2 - 2x}{-2,-1,0,1,2,3}\end{varwidth}}}{}
\end{pfloesung}
\begin{pfloesung}{}
\lz{13b)}{$1;\ 11;\ -12;\ -7$}{a) $4 + (-3) = 1$; b) $2 \cdot 4 - (-3) = 8 + 3 = 11$; c) $4 \cdot (-3) = -12$; d) $-3 - 4 = -7$}{}
\end{pfloesung}
\begin{pfloesung}{}
\lz{14.}{$-2$}{$\dfrac{1 - (-7)}{-4} = 8 : (-4) = -2$}{}
\end{pfloesung}
\begin{pfloesung}{}
\lz{15.}{$-2{,}5$}{$\dfrac{6 + (-1)}{-2} = 5 : (-2) = -2{,}5$}{}
\end{pfloesung}
\begin{pfloesung}{}
\lz{16.}{$b - a$ (Wert $6$)}{$4 - (-2) = 6$. Die anderen: $a \cdot b = (-2) \cdot 4 = -8$; $\dfrac{b}{a} = 4 : (-2) = -2$; $a - b = -2 - 4 = -6$.}{}
\end{pfloesung}
\begin{pfloesung}{}
\lz{17.}{$14$}{$6 - 2 \cdot (-4) = 6 - (-8) = 6 + 8 = 14$}{}
\end{pfloesung}
\begin{pfloesung}{}
\lz{18.}{$-20$}{$4 \cdot (-3 - 2) = 4 \cdot (-5) = -20$}{}
\end{pfloesung}
\begin{pfloesung}{}
\lz{19.}{$15$}{$2 \cdot (-3)^2 + (-3) = 2 \cdot 9 - 3 = 18 - 3 = 15$}{}
\end{pfloesung}
\begin{pfloesung}{}
\lz{20.}{$1{,}5$}{$\dfrac{-1 - 5}{-4} = (-6) : (-4) = 1{,}5$}{}
\end{pfloesung}
\begin{pfloesung}{}
\lz{21.}{$14$}{$(-3)^2 - (-5) = 9 + 5 = 14$}{}
\end{pfloesung}
\begin{pfloesung}{}
\lz{22.}{a) $18$\,m, ja \quad b) $40$\,m, nein}{a) $\dfrac{30}{10} = 3$; $3 \cdot 3 + 3^2 = 9 + 9 = 18$\,m; $18 < 20$: ja, es hält 2\,m vor dem Kind. b) $\dfrac{50}{10} = 5$; $3 \cdot 5 + 5^2 = 15 + 25 = 40$\,m; $40 > 20$: nein, der Weg ist 20\,m zu lang.}{}
\end{pfloesung}
\end{document}
